% =====================================================================
%  AI4Good workshop variant.
%  Defines everything main.tex pulls in: title, venue string, abstract,
%  introduction, discussion. Nothing else in the paper changes.
% =====================================================================

\newcommand{\wkTitle}{Portable Causal Fairness: Graph Surgery for Private
Synthetic Data \st{I'm just concerned that, despite "fairness" in the first clause, that it might signal privacy on synthetic as the key focus of this paper, which I would think is less applicable to AI4Good, and seems already solved.}}

\newcommand{\wkVenue}{AI for Good}

% ---------------------------------------------------------------- abstract
\newcommand{\wkAbstract}{%
Releasing synthetic data in place of real records gives an institution a lever it
never had over the real thing: it chooses the process that generates the table, and
can build that process so the unfair pathways are absent. DECAF made this concrete
by wiring a GAN's generator to a causal graph and severing the edges that carry
disparity from a protected attribute to an outcome, turning three fairness
definitions into three sets of edge cuts. The mechanism has never left that one
architecture, which is a practical problem: an organisation cannot adopt it without
adopting DECAF's GAN, and that GAN offers neither a privacy guarantee nor
competitive data quality. Moving it is not a matter of swapping backbones. DECAF cuts edges in a graph
somebody handed it, after training --- pure post-processing. The differentially
private synthesizers that dominate real tabular releases have no supplied graph;
they must buy one from the privacy budget, so the same definition has to be
enforced \emph{inside} the private structure search. \st{That last DP sentence is perhaps too much detail for the abstract, and not necessarily true (RAP for example, or private-pgm are given graphs. Please stop overstating things. We can sit back a little more.} These are two different
operations meeting one definition. We implement both and attach all three
definitions to nine generators from three families (marginals-based, GAN-based and
diffusion-based, each with differentially private variants) over 2{,}520
matched-pair runs. The definitions transfer: cutting edges reduces downstream
disparity in every family, dataset and privacy level, on Adult by $0.055$ of parity
gap  \st{this is a bad, unintroduced, unimpactful metric for the abstract. Just say something for the result like "we show that the fairness mechanism works well across all generation paradigms, and has <X> general consequences for privacy and <Y> general consequences for utility." That's it. } for under one accuracy point, with no loss at the strongest privacy level we
test. How much they remove does not, and the cause is structural: a cut can only
remove disparity the generator's graph encoded, and a graph bought under noise
encodes less of it \st{remove this. It's just too abstract and specific and nuanced and unsubstantiated for the paper, much less the abstract.}. On Adult, the cut removes $89\%$ of a causal diffusion
backbone's parity gap but $29\%$ of a marginals synthesizer's. That difference
decides where the intervention is worth placing. It is also the only fairness
intervention that ships with the released table, and so the only one the publisher
can guarantee. \st{this doesn't make sense. Please just simplify the abstract to not have these fabulous, unsubstantiated, speculative details.}%
}

% ------------------------------------------------------------------- intro
\newcommand{\wkIntro}{%
Institutions holding sensitive records increasingly publish a synthetic substitute
rather than the records themselves: a hospital, a lending regulator or a
statistical agency fits a generative model and releases a table drawn from it.
That substitution is usually justified on privacy grounds, but it carries a
second, less-exploited consequence. Whoever generates the data chooses the
generating process, and a process can be edited. A disparity that exists in the
world because of a pathway running from a protected attribute, such as national
origin, to an outcome, such as whether a loan is approved, need not exist in the
released table. \st{this is a good paragraph. thank you.}

This is an unusually good place to intervene. Downstream fairness corrections have
to be applied again by every analyst who touches the release, and a public table
has consumers its publisher will never meet; an edit to the generating process
ships with the data \st{didn't find this sentence clear or necessarily true}. But the one mechanism that performs such an edit, DECAF
\citep{vanbreugel2021decaf}, arrives welded to a GAN that offers no privacy
guarantee, which is the opposite of what institutions in this position actually
deploy. An agency already running a differentially private marginals-based
pipeline cannot adopt causal fairness without replacing its infrastructure with an
architecture that would degrade both its data quality and its privacy posture. \st{can perhaps mention the high utility of diffusion briefly here, since the intro should be broad.} The capability and the deployment setting have never been in the same system. \st{not sure what this sentence means.}

We ask three questions. Is edge surgery a property of DECAF, or of causal
factorisation itself --- does the mechanism still work when the generator is not a
GAN? If a synthesizer is not handed a causal graph but must infer one under a
privacy budget, what does the same fairness definition even mean, and what does it
still accomplish \st{make this second question more general, and less making the unnecessary point about how marginals-based have to do something differently to get the graph, like "can this causal graph fairness mechanism be applied to undirected graphs" or "can any generator with features related via a graph have this fairness mechanism applied to it, (directed or undirected)"}? And how does an intervention placed inside the generator compare
with the established practice of correcting fairness downstream, on identical data?
\S\ref{sec:related} places these against what exists, \S\ref{sec:method} states
what we built, and \S\ref{sec:results} answers them.%
}

% -------------------------------------------------------------- discussion
\newcommand{\wkDiscussion}{%
\textbf{The cut is only as strong as the graph it is given.} DECAF severs edges in
a causal DAG \st{I know DAG is a common acronym among some, but for the AI4good venue, can we spell it out in the intro when we first mention it?} that somebody supplied, at sampling time  \st{are you doubly-sure that DECAF doesn't have some mechanism to choose / build the graph?}. A differentially private
marginals-based synthesizer is supplied with nothing: it buys a structure out of
the privacy budget, so the same definition has to be enforced during structure
selection, against whatever graph the noise could afford. That is why the
definitions port cleanly while what they accomplish does not, and the limiting
case is stark \st{this whole sentence is weirdly imprecise and confusing. I would think the definitions are ported differently depending on the way the generator engages with / builds a graph, but the definition is similar.}. Fairness through unawareness (FTU) on MST is a literal no-op in all $45$ matched
Adult pairs, because the direct protected-to-income edge \st{I'm saying this again: we need to offer more clearly what the features are classified as and what the graph connections tend to be, so that we can see this in action.} is one the private
structure search never buys; the constraint is satisfied before it is imposed.
This reframes what an audit should ask for. Knowing which generator produced a
release tells you little. Knowing the graph it sampled along, and which attributes
were nominated admissible, tells you almost everything \st{I don't think our results substantiate this... geeze. Please be more careful, and precise and less vague.}, and on Adult moving
between our three role splits changes what conditional fairness accomplishes by
$0.033$, against $0.002$ for moving between architectures \st{what is this sentence??}. The decisive choice is
the one a person makes about which explanations are legitimate, not the one an
engineer makes about an objective \st{what??}.

\textbf{A cut inside the generator competes with downstream mitigation, and is the
only one the publisher can guarantee.} We ran an existing benchmark of eight
fairness mitigators \citep{angelozzi2026where} on the same synthesizer, splits,
privacy levels, classifier and metric, and added our in-generator cut as a stage it
does not have. Where the generator encoded the disparity in the first place, the
cut is the strongest intervention available: on Adult at $\varepsilon{=}1$ it
removes $89\%$ of the parity gap, more than any of the eight that leaves a usable
model behind --- calibrated equalised odds, the nearest such competitor, removes
$71\%$ and keeps $2.6$ more points of accuracy. (Two of the eight score better on
paper only by destroying the model, which is the subject of the next point.) Where the generator did not encode it --- COMPAS, where the privately
selected structure barely carries the protected-to-outcome dependence --- there is
little to sever and downstream mitigators do better, exactly as the argument above
predicts. So the honest summary is that the two are complements whose relative
value depends on what the graph captured. What is not symmetric is enforceability:
a downstream mitigator is applied by whoever trains a model, and a released table
is trained on by analysts the publisher will never meet, whereas a cut inside the
generator ships with the table. Fairness that depends on every future consumer
doing the right thing is not a property of the release.

\textbf{A fairness number without a utility number beside it is unfalsifiable.} \st{This conclusion feels a little weak and uninteresting + unsubstantiated... It's perhaps too niche to be of interest. Anything else revealing / impactful / novel beyond what DECAF shows?}
Every metric here rewards a table for containing no disparity, and a table that
teaches a classifier nothing contains none. This is not a hypothetical: in the
benchmark above, learned fair representations score a perfect parity gap of
exactly $0.000$ on Adult by predicting a single class for every row --- zero
recall, accuracy equal to the majority-class rate --- at both privacy levels, and
the fairness metric alone cannot tell that apart from a genuine success. Nor is the
trap ours to point at from a safe distance: our own most fairness-favourable cell
falls into it too, which is how we found it. Reporting utility beside every
fairness number is what makes the fairness number mean anything, and it costs one
extra evaluation on data the publisher already holds.%
}
